\documentclass[journal,onecolumn]{IEEEtran}
\usepackage{amsmath,amsfonts}
\usepackage{algorithmic}
\usepackage{algorithm}
\usepackage{array}
\usepackage[caption=false,font=normalsize,labelfont=sf,textfont=sf]{subfig}
\usepackage{textcomp}
\usepackage{stfloats}
\usepackage{url}
\usepackage{verbatim}
\usepackage{graphicx}
\usepackage{cite}
\hyphenation{op-tical net-works semi-conduc-tor IEEE-Xplore}
% updated with editorial comments 8/9/2021

\begin{document}

\title{Two-page Coursework for \emph{RANC}}

\author{
Name: San Zhang\\
Stu. No. 2022109
}

\maketitle

\begin{abstract}
Please manually place the \emph{IEEETran.cls} file in the tex folder if compilation does not work well.
\end{abstract}

\section{Introduction}
This is a coursework template for \emph{Recent Advances in Network Computing}. Requirements are listed as follows. 
\begin{itemize}
    \item Combine your research topic with network computing, in any aspects that you think make sense. 
    \item To ensure review quality, the length should be no more than two pages.
    \item Only \emph{necessary} references are required to be cited (eight or fewer references recommended).
    \item Only one figure is allowed, to illustrate the motivation or basic idea of this coursework.
\end{itemize}

\section{Motivation and basic idea}
In this section, you need to address one key problem:
How network computing impacts your target problem, or how your research tools can be used to enhance network computing.
\section{Key challenges}

\section{Potential solutions/Discussion}

\section{Conclusion}

\end{document}


